\documentclass[11pt]{article}

\usepackage[T1]{fontenc}
\usepackage[utf8]{inputenc}
\usepackage{amsmath,amssymb,amsthm}
\usepackage{booktabs}
\usepackage{float}
\usepackage{geometry}
\usepackage{hyperref}
\usepackage{url}

\geometry{margin=1in}
\emergencystretch=3em

\newtheorem{theorem}{Theorem}
\newtheorem{computationalresult}{Computational Result}
\newtheorem{proposition}{Proposition}
\newtheorem{lemma}{Lemma}
\newtheorem{corollary}{Corollary}
\newtheorem{conjecture}{Conjecture}
\newtheorem{openproblem}{Open Problem}
\theoremstyle{definition}
\newtheorem{definition}{Definition}
\newtheorem{example}{Example}
\theoremstyle{remark}
\newtheorem{remark}{Remark}

\numberwithin{equation}{section}

\newcommand{\Z}{\mathbb{Z}}
\newcommand{\Zp}{\mathbb{Z}^{+}}
\newcommand{\vq}[1]{v_{#1}}

\title{Computational Verification of Sun's Conjecture~4.3(i):\\
Exceptional Structure of OEIS Sequence~A248044}
\author{Carlo Corti\\
Independent researcher\\
\texttt{carlo.corti@outlook.com}}
\date{March 2026}

\begin{document}

\maketitle

\begin{abstract}
We present the first large-scale computational investigation of
OEIS (On-Line Encyclopedia of Integer Sequences) sequence A248044,
which encodes Sun's Conjecture~4.3(i):
for every positive integer $m$, there exists a positive integer~$n$
such that $(m+n) \mid \pi(m)^2 + \pi(n)^2$, where $\pi(x)$ denotes
the prime-counting function.
We determine $a(m) = \min\{n \in \Zp \text{ (positive integers)} : (m+n) \mid \pi(m)^2 + \pi(n)^2\}$
for $99\,998$ of the first $100\,000$ positive integers, with a
final search bound of $n \le 10^{13}$.
The computation employs a segmented sieve of Eratosthenes coupled
with a modular obstruction filter based on quadratic residuacity of~$-1$,
parallelised via OpenMP on~$16$ threads.
The two unresolved values, $m = 19\,623$ and $m = 19\,624$, satisfy
$a(m) > 10^{13}$, yielding Cram\'er-type ratios
$a(m)/m > 5.095 \times 10^8$ --- the largest lower bounds in the
entire sequence.
We prove that the obstruction filter is sound
(Obstruction Filter Lemma) and analyse the arithmetic
structure of the unresolved cases: all belong to a cluster sharing
$\pi(m) = 2225 = 5^2 \times 89$, whose prime factorisation consists
entirely of primes congruent to $1 \pmod{4}$, so that the
implemented obstruction filter operates at full scope.
Statistical analysis reveals a heavy-tailed distribution
consistent with a log-exponential model, and we provide
probabilistic estimates for the unresolved values.
\end{abstract}

\noindent\textbf{2020 Mathematics Subject Classification.}
Primary 11A41; Secondary 11Y55, 11N05, 11A07.

\noindent\textbf{Key words and phrases.}
prime counting function, Sum of two squares,
Computational verification, OEIS A248044, Sun's conjecture,
Obstruction filter.

\section{Introduction}
\label{sec:intro}

In a 2013 preprint subsequently published in 2019,
Zhi-Wei Sun~\cite{Sun2019} proposed a collection of conjectures
involving divisibility conditions on prime-related functions.
Among these, Conjecture~4.3(i) (p.~15 of~\cite{Sun2019}) states:

\begin{conjecture}[Sun, 2013/2019]\label{conj:sun43i}
For every positive integer $m$, there exists a positive integer $n$
such that
\begin{equation}\label{eq:divisibility}
(m + n) \mid \pi(m)^2 + \pi(n)^2,
\end{equation}
where $\pi(x) = \#\{p \le x : p \text{ prime}\}$ denotes the
prime-counting function.
\end{conjecture}

The sequence $a(m)$, defined as the least positive integer~$n$
satisfying~\eqref{eq:divisibility}, was submitted to the OEIS in
2014 and appears as sequence A248044~\cite{OEIS_A248044}.
Prior to this work, the OEIS b-file contained values only for
$m \le 200$.
To the best of our knowledge, no large-scale computational
verification of Conjecture~\ref{conj:sun43i} has appeared in the
literature (for a related computational study of a different Sun
conjecture, see~\cite{MacTsop2021}).

The companion conjecture, 4.1(i), involving $p_m$ (the $m$-th prime)
in place of $\pi(m)$, was investigated computationally
in~\cite{CortiA247975}; the present paper extends the same
methodology to the prime-counting variant.

The present work makes three main contributions.

\emph{First}, we extend the computation of $a(m)$ to
$m = 1, \ldots, 100\,000$, resolving $99\,998$ of $100\,000$ cases
with a final search bound of $n \le 10^{13}$.  This is the largest
verified dataset for this conjecture.

\emph{Second}, we carry out a systematic statistical analysis of the
distribution of $a(m)$, revealing a heavy-tailed structure well
described by a log-exponential model.

\emph{Third}, we identify and characterise the two values of~$m$ for
which the conjecture resists verification beyond $n = 10^{13}$,
both belonging to an arithmetically structured cluster sharing a
common value $\pi(m) = 2225 = 5^2 \times 89$; we explain how the
factorisation of $\pi(m)$ into primes $\equiv 1 \pmod{4}$ ensures
that the implemented obstruction filter operates at full scope,
and propose the anomalous resistance of
these two cases as an open problem.

The paper is organised as follows.
Section~\ref{sec:background} reformulates the divisibility condition
in modular arithmetic and establishes the theoretical basis for the
obstruction filter.
Section~\ref{sec:algorithm} describes the algorithm and implementation.
Section~\ref{sec:results} presents the computational results.
Section~\ref{sec:statistics} carries out the statistical analysis.
Section~\ref{sec:exceptional} analyses the two unresolved cases.
Section~\ref{sec:conclusions} concludes with open problems.

\section{Mathematical Background}
\label{sec:background}

For positive integers $m$ and $n$, set $s = m + n$.
Throughout this section, $\vq{q}(s)$ denotes the $q$-adic
valuation of~$s$, i.e., the largest exponent~$e$ such that
$q^e \mid s$.
The divisibility condition~\eqref{eq:divisibility} is equivalent to
\begin{equation}\label{eq:modular}
\pi(n)^2 \equiv -\pi(m)^2 \pmod{s}.
\end{equation}
When $\gcd(\pi(m), s) = 1$, this simplifies to
\begin{equation}\label{eq:qr}
\bigl(\pi(n) \cdot \pi(m)^{-1}\bigr)^2 \equiv -1 \pmod{s},
\end{equation}
so a solution exists at index~$n$ only if $-1$ is a quadratic
residue modulo~$s$ (a necessary but not sufficient condition, since
$\pi(n)$ must also fall in the correct residue class).

A classical result (see \cite{IrelandRosen}, Ch.~5) gives:

\begin{proposition}\label{prop:qr}
$-1$ is a quadratic residue modulo $s > 1$ if and only if both
of the following hold:
\begin{enumerate}
\item[\textup{(a)}] $4 \nmid s$
\textup{(}equivalently, $\vq{2}(s) \le 1$,
where $\vq{2}(s)$ denotes the $2$-adic valuation of~$s$\textup{)};
\item[\textup{(b)}] no prime $q \equiv 3 \pmod{4}$ divides~$s$.
\end{enumerate}
\end{proposition}

\begin{remark}\label{rem:jacobi}
Condition~(b) is strictly stronger than requiring the Jacobi
symbol $\bigl(\frac{-1}{s}\bigr) = 1$.
For odd~$s$, the Jacobi symbol equals~$1$ if and only if
$\sum_{q \equiv 3\,(4)} \vq{q}(s)$ is even; actual solvability of
$x^2 \equiv -1 \pmod{s}$ requires that no such prime divides~$s$
at all.
A standard example: $s = 21 = 3 \times 7$, where
$\bigl(\frac{-1}{21}\bigr) = 1$ yet $-1$ is not a QR modulo~$21$.
\end{remark}

\subsection{The obstruction filter}

We now formalise the obstruction that arises when a small prime
$q \equiv 3 \pmod{4}$ divides $s = m + n$ to an odd power but
does not divide~$\pi(m)$.

\begin{lemma}[Obstruction Filter]\label{lem:obstruction}
Let $q$ be a prime with $q \equiv 3 \pmod{4}$, and let
$s = m + n$ with $m, n \in \Zp$.
Suppose that $q \mid s$ and that $q \nmid \pi(m)$.
Then $s \nmid \pi(m)^2 + \pi(n)^2$.
\end{lemma}

\begin{proof}
Suppose for contradiction that $s \mid \pi(m)^2 + \pi(n)^2$.
Since $q \mid s$, we have $q \mid \pi(m)^2 + \pi(n)^2$.

Since $q \nmid \pi(m)$ by hypothesis, we have
$\gcd(\pi(m), q) = 1$, so $\pi(m)$ is invertible modulo~$q$.
From $q \mid \pi(m)^2 + \pi(n)^2$ we get
\[
\bigl(\pi(n) \cdot \pi(m)^{-1}\bigr)^2 \equiv -1 \pmod{q}.
\]
Since $q \equiv 3 \pmod{4}$, the element $-1$ is not a quadratic
residue modulo~$q$ (\cite{IrelandRosen}, Ch.~5, Theorem~3).
This contradicts the equation above, so the assumption
$s \mid \pi(m)^2 + \pi(n)^2$ must be false.
\end{proof}

\begin{corollary}\label{cor:maxobstruction}
If every prime factor of $\pi(m)$ is congruent to $1 \pmod{4}$
\textup{(}i.e., $\pi(m)$ has no prime factor $q \equiv 3 \pmod{4}$\textup{)},
then for any prime $q \equiv 3 \pmod{4}$ with $q \mid s$,
the condition of
Lemma~\textup{\ref{lem:obstruction}} is satisfied and~$s$ is
blocked.
In this case the coprimality condition $q \nmid \pi(m)$ holds
for every filter prime $q \in \{3, 7, 11, 19\}$, so the
implemented filter operates at its full scope.
\end{corollary}

\begin{proof}
If every prime factor of $\pi(m)$ is $\equiv 1 \pmod{4}$, then
$q \nmid \pi(m)$ for every $q \equiv 3 \pmod{4}$, and
Lemma~\ref{lem:obstruction} applies whenever $q \mid s$.
\end{proof}

\begin{remark}\label{rem:parity}
The condition $4 \mid s$ from Proposition~\textup{\ref{prop:qr}(a)}
can be analysed via a direct parity argument,
independent of the quadratic residue framework.
If $4 \mid s$ and $s \mid \pi(m)^2 + \pi(n)^2$, then
$4 \mid \pi(m)^2 + \pi(n)^2$.
Since $x^2 \bmod 4 \in \{0, 1\}$ for all integers~$x$, the sum
$\pi(m)^2 + \pi(n)^2 \bmod 4$ lies in $\{0, 1, 2\}$; it equals~$0$
if and only if both $\pi(m)$ and $\pi(n)$ are even.
Therefore, when $\pi(m)$ is odd --- as is the case for the
unresolved values $m = 19\,623$ and $m = 19\,624$, where
$\pi(m) = 2225$ --- the condition $4 \mid s$ provably precludes a
solution, regardless of whether $\gcd(\pi(m), s) = 1$.
For $m$ with $\pi(m)$ even, solutions with $4 \mid s$ are
theoretically possible.
Our implementation does not reject candidates based on
$4 \mid s$; the analysis above is provided as theoretical context
for the exceptional cases.
\end{remark}

\begin{remark}\label{rem:filter_general}
The proof of Lemma~\ref{lem:obstruction} is valid for any prime
$q \equiv 3 \pmod{4}$, not only for the small primes used in our
implementation.  Our code tests $q \in \{3, 7, 11, 19\}$ for
computational efficiency; testing all primes $q \equiv 3 \pmod{4}$
dividing~$s$ to odd power would make the odd-valuation filter
exhaustive, but at prohibitive computational cost.
\end{remark}

\begin{remark}\label{rem:implemented_filter}
The implementation tests the condition ``$\vq{q}(s)$ is odd''
rather than the weaker condition ``$q \mid s$'' stated in
Lemma~\textup{\ref{lem:obstruction}}.
Since $\vq{q}(s)$ odd implies $q \mid s$, the implemented
condition is strictly stronger: the code rejects a
\emph{subset} of the candidates that the Lemma permits
rejecting.  Candidates with even valuation $\vq{q}(s) \ge 2$
are not filtered and proceed to the full divisibility test;
no valid solution is ever missed.  In other words, the
implemented filter is \emph{sound} but not maximal with
respect to Lemma~\textup{\ref{lem:obstruction}}.
The choice of the odd-valuation test was driven by
computational efficiency: it enables a fast early-exit test
(\texttt{if(s\%q\^{}2 != 0) return 1}) that avoids
the full valuation loop in the common case $\vq{q}(s) = 1$.
\end{remark}

\subsection{Density of obstructed candidates}

Three progressively weaker obstruction conditions are relevant.
\emph{Condition~\textup{(b)}} of Proposition~\ref{prop:qr}
requires that no prime $q \equiv 3 \pmod{4}$ divides~$s$ at all.
The classical $B_1$ class (integers representable as a sum of two
squares) imposes the weaker requirement that every such prime
appears to an \emph{even} exponent; the set satisfying condition~(b)
is a proper subset of~$B_1$.
Finally, our implemented filter (Remark~\ref{rem:implemented_filter})
tests only whether $\vq{q}(s)$ is odd for a finite set of primes,
rejecting a subset of the candidates that fail the $B_1$ condition.

The Landau--Ramanujan theorem~\cite{Landau1908} gives the count of
$B_1$ integers up to~$x$ as $\sim K x / \sqrt{\ln x}$ for a
constant $K \approx 0.7642$; the asymptotic estimates rely on
Mertens-type products~\cite{Mertens1874,Tenenbaum1995}.
The count of integers satisfying the
strictly stronger condition~(b) has the same asymptotic order
$\sim C_0 x / \sqrt{\ln x}$ but with a smaller
constant~\cite{Tenenbaum1995}.
At the finite scales relevant to our computation ($x \sim 10^{9}$
to $10^{13}$), the fraction of candidates whose $q$-adic valuations
are all even for every $q \equiv 3 \pmod{4}$ (the $B_1$ condition)
is roughly $17$--$20\%$; the fraction satisfying the full
condition~(b) is smaller still.
In either case, $80\%$ or more of candidates are obstructed in
principle by odd-valuation primes.

Our implementation checks only $q \in \{3, 7, 11, 19\}$,
which empirically eliminates approximately $57\%$ of candidates
(measured across the monolithic run for $m = 1, \ldots, 100\,000$),
well below the theoretical maximum.  The gap is due to the
limited trial-division bound: candidates whose only factor
$q \equiv 3 \pmod{4}$ exceeds~$19$ are not eliminated.

\section{Algorithm and Implementation}
\label{sec:algorithm}

\subsection{Overview}

The computation proceeded in two phases:
a \emph{monolithic} phase computing $a(m)$ for all
$m = 1, \ldots, 100\,000$ with search bound $n \le 6 \times 10^9$,
followed by a series of \emph{targeted segmented} searches
extending the bound to $n = 10^{13}$ for the unresolved cases.
Both phases use a sieve of Eratosthenes~\cite{CrandallPomerance}
for prime enumeration.

All computations were performed on a personal workstation:
AMD Ryzen~9 7940HS (16 threads), 64\,GB DDR5 RAM,
running Linux Mint~22.3.  Compilation used GCC with
\texttt{-O3 -march=znver4 -fopenmp}.

\subsection[Monolithic search: sun-A248044-v5.c]{Monolithic search: \texttt{sun\_A248044\_v5.c}}

For the initial phase, the program builds a global table of
$\pi(n)$ for all $n \le N_{\max}$, stored as \texttt{uint32\_t}.
The sieve of Eratosthenes marks composites in a
\texttt{uint8\_t} array; a prefix sum then fills the $\pi$-table.
Memory consumption is $5 \cdot N_{\max}$ bytes
($\approx 30$\,GB for $N_{\max} = 6 \times 10^9$).

For each $m$ and each candidate $n$, the program sets
$s = m + n$ and checks:
\begin{enumerate}
\item \textbf{Obstruction filter:}
  for each $q \in \{3, 7, 11, 19\}$, if $\vq{q}(s)$ is odd
  and $q \nmid \pi(m)$, reject~$n$
  (see Remark~\ref{rem:implemented_filter}).
  Additionally, if $s$ is a prime with $s \equiv 3 \pmod{4}$
  and $s > \pi(m)$ (which holds for all $m \ge 2$), reject~$n$
  (since $\vq{s}(s) = 1$ is odd and $s \nmid \pi(m)$,
  Lemma~\ref{lem:obstruction} applies).
\item \textbf{Divisibility test:}
  compute $\pi(m)^2 + \pi(n)^2$ and check whether $s$ divides the
  sum.
\end{enumerate}

For $n \le 6 \times 10^9$, we have $\pi(n) \le \pi(6 \times 10^9)
\approx 2.80 \times 10^8$ and $\pi(m) \le \pi(100\,000)
= 9\,592$, so $\pi(m)^2 + \pi(n)^2 < 7.82 \times 10^{16}$, which
fits in \texttt{int64\_t} (max $\approx 9.2 \times 10^{18}$).

The monolithic phase resolved $99\,935$ of $100\,000$ cases in
approximately $10$~minutes of wall-clock time, leaving $65$
unresolved cases (the ``exceptional'' values of~$m$).

\subsection[Targeted segmented search: sun-A248044-targeted-v3.c]{Targeted segmented search: \texttt{sun\_A248044\_targeted\_v3.c}}
\label{sec:targeted}

The monolithic architecture cannot scale beyond
$\sim 7 \times 10^9$ due to the RAM ceiling of the global
$\pi$-table.
For the extended searches we developed a segmented approach:
the range $[n_{\min}, n_{\max})$ is divided into segments of width
$W = 10^9$, and for each segment $[l, l+W)$:

\begin{enumerate}
\item A segmented sieve of Eratosthenes marks composites in
  a \texttt{uint8\_t} array \texttt{seg[]} of size~$W$
  ($\approx 1$\,GB).
\item A prefix sum computes $\pi_{\mathrm{seg}}[k]
  = \#\{\text{primes in } [l, l+k]\}$ as \texttt{uint32\_t}
  ($\approx 4$\,GB), and the absolute count is recovered as
  $\pi(l+k) = \pi_{\mathrm{base}} + \pi_{\mathrm{seg}}[k]$,
  where $\pi_{\mathrm{base}} = \pi(l)$ is maintained across
  segments.
\item The inner loop iterates over all $k \in [0, W)$ in parallel
  (OpenMP, \texttt{schedule(guided)}), and for each still-unresolved
  target~$m$, applies the obstruction filter and divisibility test.
\end{enumerate}

Total memory per segment: $\approx 5$\,GB.  Since $\pi_{\mathrm{seg}}[k]$
counts primes within a single segment of width $10^9$, the offset
cannot exceed $\sim 4 \times 10^7$ and comfortably fits in
\texttt{uint32\_t}.

For $n > \sim 3.5 \times 10^9$ we have $\pi(n) > 1.6 \times 10^8$,
so $\pi(n)^2 > 2.6 \times 10^{16}$; and at the extremes of Ext-5
($n \sim 10^{13}$) we have $\pi(n) \approx 3.5 \times 10^{11}$, giving
$\pi(n)^2 \approx 1.2 \times 10^{23}$, which far exceeds the range
of \texttt{int64\_t}.  Accordingly, the segmented tool uses
\texttt{\_\_int128} (a GCC extension) uniformly for the sum
$\pi(m)^2 + \pi(n)^2$ and the modular reduction.
On x86-64, GCC compiles $\mathtt{\_\_int128} \mathbin{\%}
\mathtt{int64\_t}$ to a native \texttt{DIV} instruction
(\texttt{rdx:rax}), so the overhead relative to 64-bit arithmetic
is negligible (a few cycles per division).

\subsection{Parallelism and correctness}

The inner loop is parallelised with
\texttt{\#pragma omp parallel for schedule(guided)}~\cite{OpenMP}.
The choice of \texttt{guided} scheduling reflects the non-uniform
cost per iteration: the obstruction filter rejects $\sim 57\%$ of
candidates without performing the expensive 128-bit division,
creating significant load imbalance under static scheduling.

Target data (arrays of $m$, $\pi(m)^2$, and filter bitmasks)
are copied into Structure-of-Arrays (SoA) layout before the
parallel region.  With at most $65$ targets at $\le 24$ bytes each,
the working set of $\approx 1.5$\,kB fits entirely in the L1 cache
(32\,kB per core), ensuring that the hot loop is memory-bound on
the segment arrays rather than on target metadata.

The variable \texttt{found\_n[ai]} (recording the best solution
found so far for target~$ai$) is read without synchronisation
inside the parallel loop and written only under
\texttt{\#pragma omp critical}.  This constitutes a formal data
race under the C11/OpenMP memory model; however, on x86-64 under
the Total Store Order (TSO) memory model, an aligned 64-bit read
cannot produce a torn value.  The worst case is a stale read, which
may cause one redundant iteration but never affects correctness:
the critical section always compares and retains the global minimum.
An \texttt{omp atomic read} directive would eliminate the formal
race but at a cost of $\sim 4$ cycles $\times 6.5 \times 10^{10}$
iterations $\approx 87$~seconds per segment --- unacceptable given
that the benign race preserves correctness on the target
architecture.
This correctness argument is specific to the x86-64 TSO memory
model and would not hold on architectures with weaker memory
ordering (e.g., ARM); however, all computations for this paper
were performed on x86-64.

\subsection{Code review}

The targeted search code (v3) underwent two rounds of independent
review with assistance from AI language models (see the AI~Disclosure
in Section~\ref{sec:ai-disclosure}).
The first round identified two critical bugs
(integer overflow in $\pi(n)^2$ and a race condition on the
minimum-tracking logic); both were fixed in v2.
The second round confirmed correctness and led to
minor improvements (input parsing hardening, overflow fix in the
CSV output column).
All observations and their dispositions are documented in the source
code header.

\section{Computational Results}
\label{sec:results}

\subsection{Campaign summary}

Table~\ref{tab:campaigns} summarises the progressive resolution of
the $65$ initially unresolved cases across the monolithic and
extended campaigns.

\begin{table}[ht]
\centering
\caption[Campaign summary]{Campaign summary.  The monolithic run (M1) used
\texttt{sun\_A248044\_v5.c}; all Ext runs used
\texttt{sun\_A248044\_targeted\_v3.c}.
``+Resolved'' denotes newly resolved cases in each run.}
\label{tab:campaigns}
\small
\begin{tabular}{llrrrl}
\toprule
Run & Bound $n$ & +Resolved & Still open & Time & Note \\
\midrule
M1 (monolithic) & $\le 6 \times 10^9$
  & $99\,935$ & 65 & $\sim$10\,min & All $m \le 10^5$ \\
Ext-2 & $6 \times 10^9 \to 10^{10}$
  & 36 & 29 & 137\,s & \\
Ext-3 & $10^{10} \to 10^{11}$
  & 20 & 9 & 1\,132\,s & \\
Ext-4 & $10^{11} \to 10^{12}$
  & 4 & 5 & 9\,452\,s & \\
Ext-5 & $10^{12} \to 10^{13}$
  & 3 & \textbf{2} & 86\,098\,s & \\
\bottomrule
\end{tabular}
\end{table}

The total computation time for the extended campaigns was
approximately $96\,819$ seconds ($\approx 26.9$ hours).
Including the monolithic phase ($\approx 0.17$~hours), the entire
computation required approximately $27$~hours of wall-clock time
($\approx 432$ CPU-hours on $16$ threads).
An Ext-6 campaign ($10^{13} \to 10^{14}$) would require an
estimated $\sim 240$ hours of wall-clock time on the same hardware,
assuming similar per-segment cost (actual time may vary with
changing obstruction rates at higher bounds),
and was deemed infeasible.

\subsection{Largest values of \texorpdfstring{$a(m)$}{a(m)}}

Table~\ref{tab:top10_absolute} lists the ten largest absolute
values of $a(m)$ found in the computation.

\begin{table}[ht]
\centering
\caption{Top~10 values of $a(m)$ by absolute magnitude.}
\label{tab:top10_absolute}
\small
\begin{tabular}{rrrrl}
\toprule
$m$ & $\pi(m)$ & $a(m)$ & Ratio $a(m)/m$ & Campaign \\
\midrule
$83\,115$ & $8\,117$ & $4\,629\,736\,117\,663$ & $5.57 \times 10^7$ & Ext-5 \\
$40\,963$ & $4\,289$ & $2\,650\,225\,478\,854$ & $6.47 \times 10^7$ & Ext-5 \\
$19\,620$ & $2\,225$ & $1\,150\,315\,774\,321$ & $5.86 \times 10^7$ & Ext-5 \\
$57\,919$ & $5\,867$ & $330\,283\,186\,011$ & $5.70 \times 10^6$ & Ext-4 \\
$57\,920$ & $5\,867$ & $330\,283\,186\,010$ & $5.70 \times 10^6$ & Ext-4 \\
$45\,151$ & $4\,687$ & $229\,736\,032\,374$ & $5.09 \times 10^6$ & Ext-4 \\
$93\,381$ & $9\,019$ & $190\,261\,921\,669$ & $2.04 \times 10^6$ & Ext-4 \\
$85\,504$ & $8\,321$ & $79\,705\,931\,373$ & $9.32 \times 10^5$ & Ext-3 \\
$44\,857$ & $4\,661$ & $32\,348\,326\,705$ & $7.21 \times 10^5$ & Ext-3 \\
$44\,858$ & $4\,661$ & $32\,348\,326\,704$ & $7.21 \times 10^5$ & Ext-3 \\
\bottomrule
\end{tabular}
\end{table}

\subsection{Largest Cram\'er-type ratios}

Table~\ref{tab:top10_ratio} lists the ten largest
Cram\'er-type ratios $C(m) = a(m)/m$.

\begin{table}[ht]
\centering
\caption{Top~10 values by Cram\'er-type ratio $a(m)/m$.}
\label{tab:top10_ratio}
\small
\begin{tabular}{rrrrr}
\toprule
$m$ & $\pi(m)$ & $a(m)$ & $a(m)/m$ & Campaign \\
\midrule
$40\,963$ & $4\,289$ & $2\,650\,225\,478\,854$ & $6.47 \times 10^7$ & Ext-5 \\
$19\,620$ & $2\,225$ & $1\,150\,315\,774\,321$ & $5.86 \times 10^7$ & Ext-5 \\
$83\,115$ & $8\,117$ & $4\,629\,736\,117\,663$ & $5.57 \times 10^7$ & Ext-5 \\
$57\,919$ & $5\,867$ & $330\,283\,186\,011$ & $5.70 \times 10^6$ & Ext-4 \\
$57\,920$ & $5\,867$ & $330\,283\,186\,010$ & $5.70 \times 10^6$ & Ext-4 \\
$45\,151$ & $4\,687$ & $229\,736\,032\,374$ & $5.09 \times 10^6$ & Ext-4 \\
$93\,381$ & $9\,019$ & $190\,261\,921\,669$ & $2.04 \times 10^6$ & Ext-4 \\
$85\,504$ & $8\,321$ & $79\,705\,931\,373$ & $9.32 \times 10^5$ & Ext-3 \\
$44\,857$ & $4\,661$ & $32\,348\,326\,705$ & $7.21 \times 10^5$ & Ext-3 \\
$44\,858$ & $4\,661$ & $32\,348\,326\,704$ & $7.21 \times 10^5$ & Ext-3 \\
\bottomrule
\end{tabular}
\end{table}

\subsection{Distribution by order of magnitude}

Table~\ref{tab:distribution} shows the distribution of $a(m)$
across orders of magnitude, among the $99\,998$ resolved cases.

\begin{table}[ht]
\centering
\caption{Distribution of $a(m)$ by order of magnitude
($99\,998$ resolved cases plus $2$ unresolved).}
\label{tab:distribution}
\begin{tabular}{lrr}
\toprule
Range & Count & Cumulative \\
\midrule
$a(m) \le 10^6$     & $89\,685$ & $89\,685$ \\
$10^6 < a(m) \le 10^7$   & $7\,516$ & $97\,201$ \\
$10^7 < a(m) \le 10^8$   & $2\,009$ & $99\,210$ \\
$10^8 < a(m) \le 10^9$   & $576$ & $99\,786$ \\
$10^9 < a(m) \le 10^{10}$  & $185$ & $99\,971$ \\
$10^{10} < a(m) \le 10^{11}$ & $20$ & $99\,991$ \\
$10^{11} < a(m) \le 10^{12}$ & $4$ & $99\,995$ \\
$10^{12} < a(m) \le 10^{13}$ & $3$ & $99\,998$ \\
$a(m) > 10^{13}$ (unresolved) & $2$ & $100\,000$ \\
\bottomrule
\end{tabular}
\end{table}

The monotonically decreasing counts across orders of magnitude
confirm the heavy-tailed nature of the distribution.
The median value is $a(m) = 52\,846$.

\subsection{Independent verification of extended-campaign results}

All $63$ results from the extended campaigns (Ext-2 through Ext-5)
were independently verified in two steps.
First, the prime-counting values $\pi(a(m))$ reported in the CSV
output files were recomputed using
\texttt{primecount}~7.10~\cite{primecount}, an implementation of
the Deleglise--Rivat algorithm that is entirely independent of the
segmented sieve of Eratosthenes used in our search programs.
In all $63$ cases, the \texttt{primecount} output matched the
CSV value exactly, confirming the correctness of the sieve-based
$\pi(n)$ computation at the witness points.
Second, the divisibility condition
\begin{equation}\label{eq:verify}
(\pi(m)^2 + \pi(a(m))^2) \bmod (m + a(m)) = 0
\end{equation}
was verified directly using the \texttt{primecount}-confirmed
values and Python's arbitrary-precision integers, with zero
discrepancies.
The full verification log is available in the project
repository~\cite{GithubRepo}.

The values of $\pi(a(m))$ span from
$\approx 2.8 \times 10^8$ (for $a(m) \approx 6 \times 10^9$,
Ext-2) to $\approx 1.6 \times 10^{11}$
(for $a(m) \approx 4.6 \times 10^{12}$, Ext-5).
As an additional cross-check, Table~\ref{tab:pi_checkpoints}
lists standard checkpoints $\pi(10^k)$ from OEIS
A006880~\cite{OEIS_A006880}.

\begin{table}[ht]
\centering
\caption{Checkpoints for $\pi(n)$ used in independent verification
(OEIS A006880, confirmed with SymPy~1.13~\cite{SymPy}).}
\label{tab:pi_checkpoints}
\begin{tabular}{rr}
\toprule
$n$ & $\pi(n)$ \\
\midrule
$10^{8}$ & $5\,761\,455$ \\
$10^{9}$ & $50\,847\,534$ \\
$6 \times 10^{9}$ & $279\,545\,368$ \\
$10^{10}$ & $455\,052\,511$ \\
$10^{11}$ & $4\,118\,054\,813$ \\
$10^{12}$ & $37\,607\,912\,018$ \\
$10^{13}$ & $346\,065\,536\,839$ \\
\bottomrule
\end{tabular}
\end{table}

\subsection{Consecutive-\texorpdfstring{$m$}{m} pairs and the shadowing mechanism}

A striking feature of the dataset is the prevalence of consecutive
pairs $(m, m+1)$ sharing large values of $a$.  The top-10 tables
include the pairs $(57\,919, 57\,920)$ and $(44\,857, 44\,858)$,
each with $a(m+1) = a(m) - 1$.  The following proposition
explains this phenomenon.

\begin{proposition}[Shadowing]\label{prop:shadowing}
Let $m \ge 1$ and suppose $a(m) > 1$.
If $m + 1$ is composite \textup{(}i.e., $\pi(m+1) = \pi(m)$\textup{)}
and $a(m)$ is composite \textup{(}i.e., $\pi(a(m)) = \pi(a(m)-1)$\textup{)},
then $a(m+1) \le a(m) - 1$.
\end{proposition}

\begin{proof}
Set $n^* = a(m) - 1$.  Then $n^* \ge 1$ since $a(m) > 1$.
Since $m + 1$ is composite, $\pi(m+1) = \pi(m)$.
Since $a(m)$ is composite, $\pi(n^*) = \pi(a(m) - 1) = \pi(a(m))$.
The new modulus is
\[
s' = (m+1) + n^* = (m+1) + (a(m) - 1) = m + a(m) = s,
\]
where $s = m + a(m)$ is the modulus from the solution for~$m$.
Therefore
\[
\pi(m+1)^2 + \pi(n^*)^2 = \pi(m)^2 + \pi(a(m))^2,
\]
which is divisible by $s = s'$ by the definition of $a(m)$.
Hence $n^*$ is a valid witness for $m + 1$, giving
$a(m+1) \le n^* = a(m) - 1$.
\end{proof}

\begin{remark}
The shadowing mechanism is not guaranteed: either $m+1$ or $a(m)$
could be prime, in which case the hypotheses of
Proposition~\textup{\ref{prop:shadowing}} fail.
When both conditions hold, the proposition produces exact pairs
with $a(m+1) = a(m) - 1$ provided $n^* = a(m) - 1$ is indeed
the minimum.
Empirically, equality $a(m+1) = a(m) - 1$ holds in all observed
instances, suggesting that when the shadowing mechanism applies,
$n^* = a(m) - 1$ is typically the first valid witness for $m+1$.
\end{remark}

\section{Statistical Analysis}
\label{sec:statistics}

\subsection{Heavy-tailed distribution}

The distribution of $a(m)$ is strongly heavy-tailed
(cf.~\cite{Hill1975} for the Hill estimator framework): while the
median is approximately $52\,846$, the maximum resolved value is
$\approx 4.63 \times 10^{12}$, a ratio of $\sim 8.8 \times 10^7$.
Values of $a(m)$ exceeding $10^9$ occur for $212$ values of~$m$
out of $99\,998$ resolved ($0.212\%$), while values exceeding
$10^{12}$ occur for $3$ values ($0.003\%$).

\subsection{A log-exponential model}

We model the occurrence of solutions heuristically.
For a given~$m$, each candidate $s = m + n$ (not blocked by the
obstruction filter) independently has a ``pass probability'' that
depends on the number of valid residue classes for $\pi(n)$
modulo~$s$.  Under a Cram\'er-type heuristic~\cite{Cramer1936,Granville1995},
the number of
solutions in an interval $[N_1, N_2]$ is approximately proportional
to $\ln(N_2/N_1)$, giving a log-exponential distribution:
\begin{equation}\label{eq:logexp}
\Pr\bigl(a(m) < N \mid a(m) > N_0\bigr)
= 1 - \exp\bigl(-\lambda \cdot \ln(N/N_0)\bigr)
= 1 - (N_0/N)^{\lambda},
\end{equation}
where $\lambda > 0$ is a pass-rate parameter.

\subsection{Calibration from extended campaigns}

The pass-rate parameter~$\lambda$ can be estimated from the
extended campaigns by the maximum-likelihood relation
\begin{equation}\label{eq:mle}
\hat\lambda = \frac{-\ln(1 - r/t)}{\ln(N_{\max}/N_{\min})},
\end{equation}
where $r$ and $t$ denote the number of cases resolved and the
number of targets entering the campaign.
(The first-order Taylor approximation
$\lambda \approx r / (t \ln(N_{\max}/N_{\min}))$
underestimates~$\hat\lambda$ when $r/t$ is large.)
Table~\ref{tab:passrate} summarises the estimates.

\begin{table}[ht]
\centering
\caption{Empirical pass rates by campaign, estimated via
Equation~\eqref{eq:mle}.  The declining
trend reflects the selection effect: cases surviving to higher
bounds have intrinsically lower pass rates.}
\label{tab:passrate}
\begin{tabular}{lrrl}
\toprule
Campaign & Resolved/Targets & $\ln(N_{\max}/N_{\min})$ &
  $\hat\lambda$ \\
\midrule
Ext-2 & $36/65$  & $0.511$ & $1.58$ \\
Ext-3 & $20/29$  & $2.303$ & $0.51$ \\
Ext-4 & $4/9$    & $2.303$ & $0.26$ \\
Ext-5 & $3/5$    & $2.303$ & $0.40$ \\
\bottomrule
\end{tabular}
\end{table}

The declining pass rate from Ext-2 to the later campaigns
reflects a strong selection effect: the cases surviving to higher
bounds are precisely those with the highest intrinsic obstruction,
so the average pass rate among survivors decreases with the bound.

Using the Ext-5 estimate $\hat\lambda \approx 0.40$ for the
two unresolved cases, the conditional median (given $a(m) > 10^{13}$)
is
\[
\tilde{N} = 10^{13} \cdot \exp(\ln 2 / 0.40)
\approx 10^{13} \times 5.7
\approx 5.7 \times 10^{13}.
\]

\subsection{Probabilistic estimates for unresolved cases}

Under the log-exponential model with $\hat\lambda = 0.40$,
the \emph{heuristic} probability of resolving each unresolved case
within various bounds is:
\begin{align*}
\Pr\bigl(a(m) \le 10^{14} \mid a(m) > 10^{13}\bigr)
  &\approx 1 - 10^{-0.40} \approx 0.60, \\
\Pr\bigl(a(m) \le 10^{15} \mid a(m) > 10^{13}\bigr)
  &\approx 1 - 10^{-0.80} \approx 0.84, \\
\Pr\bigl(a(m) \le 10^{16} \mid a(m) > 10^{13}\bigr)
  &\approx 1 - 10^{-1.19} \approx 0.94.
\end{align*}

The heuristic probability of finding \emph{both} cases within
$[10^{13}, 10^{14}]$ is approximately $0.60^2 \approx 0.36$,
and within $[10^{13}, 10^{15}]$ approximately
$0.84^2 \approx 0.71$.

\begin{remark}
The log-exponential model treats $a(m)$ as though each logarithmic
unit of search space is equally likely to yield a solution.  This is
a heuristic, not a theorem; the data are deterministic, not i.i.d.,
and the model does not account for the specific arithmetic structure
of individual cases.  The estimates above should be understood as
order-of-magnitude guidance, not precise probabilities.
\end{remark}

\section{Exceptionally Resistant Cases}
\label{sec:exceptional}

\subsection{The two unresolved values}

After the full campaign (bound $n = 10^{13}$), the two values
\begin{equation}\label{eq:open}
m = 19\,623 \quad\text{and}\quad m = 19\,624
\end{equation}
remain unresolved, with
\[
a(19\,623) > 10^{13}, \qquad a(19\,624) > 10^{13},
\]
giving Cram\'er-type lower bounds
$a(m)/m > 5.095 \times 10^8$ --- the largest in the entire dataset.

The exhaustive search establishes the following.

\begin{theorem}[Computational lower bounds]\label{thm:lower}
For each $m \in \{19\,623,\; 19\,624\}$, the divisibility condition
$(m+n) \mid \pi(m)^2 + \pi(n)^2$ has no solution with
$1 \le n \le 10^{13}$.
Consequently, either $a(m) > 10^{13}$ or
Conjecture~\textup{\ref{conj:sun43i}} fails for this value of~$m$.
\end{theorem}

\begin{proof}
The full range $1 \le n \le 10^{13}$ was covered in two phases.
The monolithic program \path{sun_A248044_v5.c}
(Section~\ref{sec:algorithm}) tested every $n$ in the range
$1 \le n \le 6 \times 10^9$ for all $m \le 100\,000$,
finding no solution for the two values in~\eqref{eq:open}.
The targeted segmented program
\path{sun_A248044_targeted_v3.c}
(Section~\ref{sec:targeted}) then extended the search over
the intervals $6 \times 10^9 < n \le 10^{13}$ via the
campaigns Ext-2 through Ext-5 (Table~\ref{tab:campaigns}),
again finding no solution for these two values.
Both programs enumerate values of $\pi(n)$ via a sieve of
Eratosthenes and, for each $n$ not eliminated by the modular
filter of Section~\ref{sec:background}, evaluate the
divisibility condition
$(\pi(m)^2 + \pi(n)^2) \bmod (m+n) = 0$ directly
(using \texttt{int64\_t} in the monolithic phase and
\texttt{\_\_int128} in the segmented phase).
The modular filter is \emph{sound}
(Lemma~\ref{lem:obstruction}), including the self-blocking
prime check, which is a special case of the Lemma with $q = s$:
it eliminates only values of $n$ for
which no solution can exist, so no valid candidate is skipped.
For both values of $m$, the search returns no solution in the
stated range; the sieve is deterministic and complete, so the
conclusion follows.
All $63$ resolved extended-search cases were independently
verified: the prime-counting values $\pi(a(m))$ were recomputed
with \texttt{primecount}~\cite{primecount} and the divisibility
condition~\eqref{eq:verify} was confirmed with zero discrepancies
(Section~\ref{sec:results}).
The source code is publicly available~\cite{GithubRepo}.
\end{proof}

\begin{remark}
The result of Theorem~\textup{\ref{thm:lower}} is conditional on
the correctness of the deterministic search implementation.
All $63$ resolved extended-search cases have been independently
verified: the values $\pi(a(m))$ were recomputed using
\texttt{primecount}~\textup{\cite{primecount}}
(Deleglise--Rivat algorithm, independent of our sieve), all
$63$ values matched exactly, and the divisibility
condition~\eqref{eq:verify} was confirmed with zero
discrepancies.  This includes the three resolved cases with the
largest absolute values
\textup{(}Table~\textup{\ref{tab:top10_absolute})}.
\end{remark}

\subsection{The cluster \texorpdfstring{$\{19\,620, 19\,623, 19\,624\}$}{\{19620, 19623, 19624\}}}

The three values $m \in \{19\,620, 19\,623, 19\,624\}$ all satisfy
$\pi(m) = 2225 = 5^2 \times 89$.  The value $m = 19\,620$ was
resolved in Ext-5 at $a(19\,620) = 1\,150\,315\,774\,321$
(ratio $\approx 5.86 \times 10^7$), but $m = 19\,623$ and
$m = 19\,624$ resist beyond $10^{13}$ --- a factor of at least
approximately $8.7$ further.

The factorisation $2225 = 5^2 \times 89$ consists entirely of
primes congruent to $1 \pmod{4}$ ($5 \equiv 1 \pmod{4}$ and
$89 \equiv 1 \pmod{4}$).  By Corollary~\ref{cor:maxobstruction},
the implemented obstruction filter operates at its full scope for
all three values, since no filter prime $q \in \{3, 7, 11, 19\}$
(all $\equiv 3 \pmod{4}$) divides $\pi(m) = 2225$.

This full-scope obstruction explains why all three values are
exceptional, but it does not explain why $m = 19\,623$ and
$m = 19\,624$ resist so much more than $m = 19\,620$.

\subsection{Congruence analysis}

Table~\ref{tab:congruences} compares the residues of the three
cluster members modulo small numbers.

\begin{table}[ht]
\centering
\caption{Residues of the cluster $\{19\,620, 19\,623, 19\,624\}$
modulo small values.}
\label{tab:congruences}
\begin{tabular}{rrrrrrl}
\toprule
$m$ & $m \bmod 3$ & $m \bmod 4$ & $m \bmod 7$ & $m \bmod 11$
  & $m \bmod 19$ & Status \\
\midrule
$19\,620$ & 0 & 0 & 6 & 7 & 12 & Resolved ($1.15 \times 10^{12}$) \\
$19\,623$ & 0 & \textbf{3} & 2 & 10 & 15 & Open ($> 10^{13}$) \\
$19\,624$ & 1 & 0 & 3 & \textbf{0} & 16 & Open ($> 10^{13}$) \\
\bottomrule
\end{tabular}
\end{table}

Notable features:
\begin{itemize}
\item $m = 19\,623 \equiv 3 \pmod{4}$: for candidate
  $n \equiv 2 \pmod{4}$, $s = m + n \equiv 1 \pmod{4}$
  (favourable); but for $n \equiv 1 \pmod{4}$,
  $s \equiv 0 \pmod{4}$ (blocked by condition~(a)).
  This creates a different interaction pattern with the other
  obstruction filters than for the other two cluster members.
\item $11 \mid 19\,624$: since $11 \equiv 3 \pmod{4}$ and
  $11 \nmid \pi(19\,624) = 2225$, any $s$ with $\vq{11}(s)$ odd
  is blocked.  The constraint $11 \mid m$ forces
  $s = m + n \equiv n \pmod{11}$, creating a periodic pattern
  that may interact unfavourably with the other obstructions.
\end{itemize}

The precise mechanism by which these congruence patterns produce an
extra factor of at least approximately~$8.7$ in resistance relative to $m = 19\,620$
remains unexplained, and we propose it as an open problem
(see Section~\ref{sec:conclusions}).

\section{Conclusions and Open Problems}
\label{sec:conclusions}

We have determined $a(m)$ for $99\,998$ of the first $100\,000$
positive integers, with final search bound $n \le 10^{13}$,
providing, to the best of our knowledge, the first large-scale
computational verification of Sun's Conjecture~4.3(i).
The dataset reveals a heavy-tailed distribution
well described by a log-exponential model, with record values
reaching $a(83\,115) = 4\,629\,736\,117\,663$ and
Cram\'er-type ratios up to $6.47 \times 10^7$.

The progressive resolution of the exceptional cases --- from $65$
unresolved at bound $6 \times 10^9$, to $29$ at $10^{10}$, to~$9$
at $10^{11}$, to~$5$ at $10^{12}$, and finally to~$2$ at
$10^{13}$ --- strongly supports the conjecture: the pattern is
consistent with a heavy-tailed distribution where
every case is eventually resolved at a sufficiently high bound,
rather than with the existence of structural counterexamples.

\paragraph{Comparison with Conjecture~4.1(i).}
A companion computation~\cite{CortiA247975} applies the same
methodology to sequence A247975, which encodes
Sun's Conjecture~4.1(i): $(m+n) \mid p_m^2 + p_n^2$.
Computing $a(m)$ for $m \le 120\,000$ with bound $n \le 2 \times 10^{11}$
resolves $119\,995$ of $120\,000$ cases, leaving $5$ unresolved.
The comparison reveals instructive structural differences:
the obstruction rate is higher ($73$--$77\%$ vs.\ ${\sim}\,57\%$
here), the maximum resolved Cram\'er-type ratio is smaller
(${\sim}\,1.97 \times 10^5$ vs.\ ${\sim}\,6.5 \times 10^7$),
and no arithmetically structured cluster analogous to
$\{19\,620, 19\,623, 19\,624\}$ appears among the unresolved cases
of A247975.
These contrasts indicate that the difficulty of Sun's divisibility
conjectures depends critically on the arithmetic function involved
($p_k$ vs.\ $\pi(k)$), not only on the shared divisibility
framework $(m+n) \mid f(m)^2 + f(n)^2$.

We conclude with three open problems.

\begin{openproblem}\label{op:unresolved}
Determine $a(19\,623)$ and $a(19\,624)$.  Both satisfy
$a(m) > 10^{13}$, with $a(m)/m > 5.095 \times 10^8$.
The heuristic model suggests a median around
$5.7 \times 10^{13}$, with approximately~$84\%$ probability
of resolution within $n \le 10^{15}$.
An Ext-6 campaign ($10^{13} \to 10^{14}$) would require
$\sim 240$~hours on our hardware; algorithmic improvements
(e.g., bitset sieve with $4\times$ segment width,
extended filter with $q < 200$) might reduce this to
$\sim 35$--$40$~hours.
\end{openproblem}

\begin{openproblem}\label{op:cluster}
Explain the anomalous resistance of $m = 19\,623$ and
$m = 19\,624$ relative to $m = 19\,620$, all sharing
$\pi(m) = 2225$.  The resolved case $a(19\,620) \approx
1.15 \times 10^{12}$ is already extreme, yet the unresolved
cases exceed it by a factor of at least approximately~$8.7$.
Is there an arithmetic characterisation of the congruence
conditions on~$m$ (beyond the factorisation of~$\pi(m)$) that
predicts this extra resistance?
\end{openproblem}

\begin{openproblem}\label{op:model}
Develop a refined probabilistic model for $a(m)$ that
accounts for the dependence on the factorisation of~$\pi(m)$.
The current log-exponential model treats all unresolved cases
uniformly; a model incorporating the obstruction rate
(which depends on the prime factorisation of~$\pi(m)$ and the
congruence class of~$m$) would provide sharper predictions.
\end{openproblem}

\section*{Acknowledgements}

The author thanks Zhi-Wei Sun for the beautiful conjecture that
motivated this work.

\section*{Data availability}
\label{sec:data}

The complete dataset ($99\,998$ resolved values and the
$2$~unresolved $m$-values), the full CSV outputs for all
campaigns, and machine-readable verification logs are publicly
available at the GitHub repository
\url{https://github.com/carcorti/A248044}, the permanent Zenodo archive
\url{https://doi.org/10.5281/zenodo.19182594}, and the OEIS entry
\url{https://oeis.org/A248044}.

The C~source code for the monolithic search
(\path{sun_A248044_v5.c}), the segmented targeted search
(\path{sun_A248044_targeted_v3.c}), and a Python verification script
(\path{verify_A248044.py}) that checks the divisibility condition
$(m+a(m)) \mid \pi(m)^2 + \pi(a(m))^2$ for any supplied list
of $(m, a(m))$ pairs are available at the GitHub and Zenodo
repositories listed above.

\section*{AI use disclosure}
\label{sec:ai-disclosure}

The computational implementation, including the C~search programs
and segmented sieve, was developed with extensive assistance
from Anthropic's Claude language models
(\texttt{claude-sonnet-4-6} and \texttt{claude-opus-4-6}).
Claude generated substantial drafts of the algorithmic and
statistical sections, performed two independent code reviews
of the targeted search program, and assisted in the statistical
analysis.
All final text and numerical results were verified and approved by
the author, who bears full scientific responsibility for all results,
including those portions where AI assistance was used.

\begin{thebibliography}{99}
\sloppy
\raggedright

\bibitem{Sun2019}
Sun, Z.-W.:
Some new problems in additive combinatorics.
Nanjing Univ.\ J.\ Math.\ Biquarterly \textbf{36}(2), 134--155
(2019).
Preprint: arXiv:1309.1679 [math.NT] (submitted 2013, v9: 2020).
Conjecture~4.3(i) appears on p.~15.
\url{https://arxiv.org/abs/1309.1679}

\bibitem{CortiA247975}
Corti, C.:
Exceptional cases and statistical structure of sequence A247975:
a computational study of Sun's Conjecture~4.1(i).
Preprint, 2026; data and code: Zenodo
\url{https://doi.org/10.5281/zenodo.18920371};
OEIS b-file extension published March~2026,
\url{https://oeis.org/A247975}

\bibitem{OEIS_A248044}
OEIS Foundation:
Sequence A248044.
The On-Line Encyclopedia of Integer Sequences
(accessed March~2026).
\url{https://oeis.org/A248044}

\bibitem{OEIS_A006880}
OEIS Foundation:
Sequence A006880 ($\pi(10^n)$).
The On-Line Encyclopedia of Integer Sequences
(accessed March~2026).
\url{https://oeis.org/A006880}

\bibitem{IrelandRosen}
Ireland, K., Rosen, M.:
A Classical Introduction to Modern Number Theory,
2nd edn. Springer, New York (1990)

\bibitem{CrandallPomerance}
Crandall, R., Pomerance, C.:
Prime Numbers: A Computational Perspective,
2nd edn. Springer, New York (2005)

\bibitem{Cramer1936}
Cram\'er, H.:
On the order of magnitude of the difference between
consecutive prime numbers.
Acta Arith.\ \textbf{2}, 23--46 (1936)

\bibitem{Granville1995}
Granville, A.:
Harald Cram\'er and the distribution of prime numbers.
Scand.\ Actuarial J.\ \textbf{1995}(1), 12--28 (1995)

\bibitem{Landau1908}
Landau, E.:
\"Uber die Einteilung der positiven ganzen Zahlen in vier
Klassen nach der Mindestzahl der zu ihrer additiven Zusammensetzung
erforderlichen Quadrate.
Arch.\ Math.\ Phys.\ \textbf{13}, 305--312 (1908)

\bibitem{Mertens1874}
Mertens, F.:
Ein Beitrag zur analytischen Zahlentheorie.
J.\ Reine Angew.\ Math.\ \textbf{78}, 46--62 (1874)

\bibitem{Tenenbaum1995}
Tenenbaum, G.:
Introduction to Analytic and Probabilistic Number Theory.
Cambridge Univ.\ Press (1995)

\bibitem{Hill1975}
Hill, B.M.:
A simple general approach to inference about the tail of a
distribution.
Ann.\ Statist.\ \textbf{3}, 1163--1174 (1975).
\url{https://doi.org/10.1214/aos/1176343247}

\bibitem{SymPy}
Meurer, A., et al.:
SymPy: symbolic computing in Python.
PeerJ Computer Science \textbf{3}, e103 (2017).
\url{https://doi.org/10.7717/peerj-cs.103}

\bibitem{OpenMP}
OpenMP Architecture Review Board:
OpenMP Application Programming Interface, Version~5.2 (2021).
\url{https://www.openmp.org/specifications/}

\bibitem{MacTsop2021}
Machiavelo, A., Tsopanidis, A.:
Zhi-Wei Sun's 1-3-5 conjecture and variations.
J.\ Number Theory \textbf{222}, 135--155 (2021).
\url{https://doi.org/10.1016/j.jnt.2020.10.001}

\bibitem{GithubRepo}
Corti, C.:
A248044: computational verification of Sun's Conjecture 4.3(i).
GitHub repository and Zenodo archive (2026).
\url{https://github.com/carcorti/A248044}.
\url{https://doi.org/10.5281/zenodo.19182594}

\bibitem{primecount}
Walisch, K.:
primecount: a high-performance prime counting implementation
(Deleglise--Rivat algorithm), version~7.10 (2024).
\url{https://github.com/kimwalisch/primecount}

\end{thebibliography}

\end{document}
